\part{Les réfractaires au modèle d'entreprise libérée et leurs arguments}

Il y a deux façons de traiter ceux qui ne croient pas à votre modèle : les caricaturer, ou les écouter.

La première est confortable et ne coûte rien — c'est d'ailleurs pour ça qu'elle domine le débat, dans les deux camps. La seconde est la seule utile. Les meilleurs critiques de l'entreprise libérée ne sont pas des nostalgiques du chef à l'ancienne : ce sont des chercheurs, des praticiens et des salariés qui ont vu le modèle de près et qui pointent des failles réelles.

Cette partie leur ouvre la porte en grand. D'abord les critiques — économiques, sociales, de gouvernance — présentées sans paille ni caricature. Ensuite les alternatives : coopératives, entreprise démocratique, sociocratie, holacratie, agilité — parce qu'on ne choisit bien un modèle qu'en connaissant les autres.

Un livre de convaincu qui esquiverait cette partie ne vaudrait rien. Allons-y franchement.
